\documentclass[12pt]{article}
\usepackage[utf8]{inputenc}
\usepackage[T1]{fontenc}
\usepackage[margin=1in]{geometry}
\usepackage{setspace}
\usepackage{xcolor}
\usepackage[colorlinks=true, urlcolor=blue!60!black, linkcolor=blue!60!black]{hyperref}
\setlength{\parindent}{0em}
\setlength{\parskip}{0.45em}
\renewcommand{\baselinestretch}{1.0}
\begin{document}
\begin{center}
{\Large \textit{Lochner v. New York}, 198 U.S. 45 (1905)\par}
\vspace{0.3em}
{\normalsize Case Brief --- One-Page Edition\par}
\end{center}

\textbf{Memory Jogger.} A cap on bakers' hours violated due process ``liberty of contract'' because it lacked a direct relation to health (\href{https://caselaw.findlaw.com/court/us-supreme-court/198/45.html}{FindLaw}).

\textbf{Facts.} New York's Bakeshop Act regulated bakery sanitation and, in Section~110, barred bakery employees from working more than sixty hours a week or ten hours a day, with no emergency exception. Utica baker Joseph Lochner, a repeat offender, was indicted for permitting an employee, who wanted the extra wages, to work more than sixty hours in a week.

\textbf{Procedural History.} Convicted and fined \$50; affirmed 3--2 by the Appellate Division and 4--3 by the Court of Appeals.

\textbf{Votes.} Reversed, 5--4: Peckham for the Court; Harlan (with White and Day) and Holmes dissented.

\textbf{Holding.} Due process protects liberty to buy and sell labor; a law restricting labor contracts between competent adults must bear a direct relation to a legitimate police-power end, and Section~110 did not.

\textbf{Reasoning.} Bakers are competent adults, bread's wholesomeness does not depend on their hours, and baking is not unusually hazardous, unlike the mining in \textit{Holden v. Hardy}; the sanitary sections showed the cap was a labor law passed ``from other motives.'' Harlan saw ``room for debate'' on bakers' health and would defer. Holmes charged that the case was decided ``upon an economic theory which a large part of the country does not entertain,'' since ``[t]he Fourteenth Amendment does not enact Mr. Herbert Spencer's \textit{Social Statics}.''

\textbf{Examples.} \emph{Same-Side:} a 48-hour cap for retail clerks; a minimum wage for adult restaurant workers. \emph{Opposite-Side:} bakery ventilation and fire-exit rules; an eight-hour day in underground mines. \emph{Fence-Sitter:} a ten-hour cap for intercity bus drivers.

\textbf{Critique.} Long condemned as judicial laissez-faire and repudiated in \textit{West Coast Hotel} (1937), \textit{Lochner} is now defended by some libertarian revisionists, such as Bernstein, as principled. Sunstein faults its treatment of common-law baselines as natural and neutral; originalists like Bork reject substantive due process altogether.

\textbf{Key Quotation.} ``General propositions do not decide concrete cases. The decision will depend on a judgment or intuition more subtle than any articulate major premise.'' --- Justice Holmes, dissenting, 198 U.S. at 76.

\end{document}
